\structsection{ЗАКЛЮЧЕНИЕ}

В работе построены интерпретации основных утверждений планиметрии Лобачевского в модели Кэли~--- Клейна. Модель описана на языке абсолюта: полюс и поляра, сложное отношение, π-движения как проективные преобразования, сохраняющие форму абсолюта. Симметрия относительно прямой записана формулой~\eqref{eq:otrazhenie}, и из неё выведено правило, на котором держатся почти все построения: перпендикуляр к хорде проходит через её полюс.

С помощью этого правила в модели получены явные формулы. Если поместить вершину угла параллельности в центр круга, формула Лобачевского $\operatorname{tg}\frac{\Pi(x)}2=e^{-x}$ получается в две строки. Расстояние от точки до диаметра~\eqref{eq:do-diametra} даёт явный вид теорем о расстоянии между точками параллельных и расходящихся прямых и уравнение эквидистанты. Общий перпендикуляр расходящихся прямых~--- хорда через их полюсы, и признак расходящихся прямых сводится к тому, лежит ли точка пересечения продолжений хорд вне круга. Для четырёхугольника Саккери найдена формула угла при верхнем основании~\eqref{eq:sakkeri}. В отрицании предложения Лежандра граничная точка оказалась связана с углом параллельности: $\Pi(OP^*)$ равен данному углу.

В разделах~5 и~6 те же средства применены к пучкам и связкам прямых. Траектория пучка и поверхность, ортогональная связке, задаются одним условием $\sk{X}{V}=c$, а три типа (окружность, орицикл, эквидистанта и соответствующие поверхности) различаются только положением точки $V$: внутри круга, на абсолюте или вне его. Найдены точные изображения сферы, орисферы и эквидистантной поверхности и их внутренняя кривизна. Показано, что кубическая комната в пространстве Лобачевского не может иметь прямых углов, а комната с полом и потолком из концентрических орисфер может, причём её потолок в $e^{h}$ раз меньше пола.

По формулам работы создана веб-страница «Круг Кэли~--- Клейна» из двадцати одной сцены. В четырнадцати сценах на плоскости фигуры модели можно передвигать и сразу видеть вычисленные длины и углы, в семи сценах в пространстве показаны сфера, орисфера, эквидистантная поверхность и три комнаты. Следующие шаги~--- объединить все модели на одном сайте и построить модель комнаты, которую можно рассматривать изнутри.

\structsection{СПИСОК ИСПОЛЬЗОВАННЫХ ИСТОЧНИКОВ}

\begin{enumerate}[label=\arabic*.,leftmargin=*,labelindent=\parindent,itemsep=0pt]
  \item Атанасян, Л.\,С. Геометрия Лобачевского / Л.\,С.~Атанасян.~--- Москва\,: БИНОМ. Лаборатория знаний, 2014.~--- 320\,с.
  \item Атанасян, Л.\,С. Геометрия. Ч.\,2 / Л.\,С.~Атанасян, В.\,Т.~Базылев.~--- Москва\,: Просвещение, 1987.~--- 368\,с.
  \item Ефимов, Н.\,В. Высшая геометрия / Н.\,В.~Ефимов.~--- Москва\,: Наука, 1971.~--- 584\,с.
  \item Каган, В.\,Ф. Основания геометрии. Ч.\,2 / В.\,Ф.~Каган.~--- Москва\,; Ленинград\,: ГИТТЛ, 1949.~--- 340\,с.
  \item Норден, А.\,П. Элементарное введение в геометрию Лобачевского / А.\,П.~Норден.~--- Москва\,: ГИТТЛ, 1953.~--- 216\,с.
  \item Прасолов, В.\,В. Геометрия Лобачевского / В.\,В.~Прасолов.~--- Москва\,: МЦНМО, 2004.~--- 208\,с.
  \item Шафаревич, И.\,Р. Линейная алгебра и геометрия / И.\,Р.~Шафаревич, А.\,О.~Ремизов.~--- Москва\,: Физматлит, 2009.~--- 512\,с.
  \item Широков, П.\,А. Краткий очерк основ геометрии Лобачевского / П.\,А.~Широков.~--- Москва\,: Наука, 1983.~--- 208\,с.
  \item Ratcliffe, J.\,G. Foundations of Hyperbolic Manifolds / J.\,G.~Ratcliffe.~--- 3rd ed.~--- Cham\,: Springer, 2019.~--- (Graduate Texts in Mathematics\,; vol.~149). \todo{сверить выходные данные, число страниц}
  \item Cannon, J.\,W. Hyperbolic geometry / J.\,W.~Cannon, W.\,J.~Floyd, R.~Kenyon, W.\,R.~Parry // Flavors of Geometry / ed. by S.~Levy.~--- Cambridge\,: Cambridge University Press, 1997.~--- (MSRI Publications\,; vol.~31).~--- P.~59--115. \todo{сверить}
  \item Coxeter, H.\,S.\,M. Regular honeycombs in hyperbolic space / H.\,S.\,M.~Coxeter // Proceedings of the International Congress of Mathematicians, 1954.~--- Amsterdam, 1956.~--- Vol.~III.~--- P.~155--169. \todo{сверить}
\end{enumerate}

\structsection{ПРИЛОЖЕНИЕ А}
\begin{center}
(справочное)\\
\textbf{Компьютерные модели и численная проверка формул}
\end{center}

Интерактивная модель «Круг Кэли~--- Клейна» (плоскость и пространство): \todo{адрес страницы после публикации на сайте}.

Модели на языке Python опубликованы в репозиториях:
\begin{itemize}
  \item гиперболическая сфера: \url{https://github.com/mvbulgakova/hyperbolic-geometry-app};
  \item орисфера: \url{https://github.com/mvbulgakova/hyperbolic-geometry2-app}.
\end{itemize}

Формулы п.~5.6 проверялись программой, приведённой ниже. Для точек построенной сферы вычисляется расстояние до центра по предложению~\ref{prop:ch}, для точек орисферы~--- величина $(1-k\cdot\omega)/\sqrt{1-|k|^2}$, которая должна быть постоянной. На 200 случайных наборах параметров (центр на расстоянии до $0{,}95$ от центра шара, $0{,}05\le\rho\le2{,}5$, $0{,}05\le r\le0{,}95$) наибольшее отклонение составило $4\cdot10^{-13}$ для сферы и $4\cdot10^{-10}$ для орисферы; для центров на расстоянии $0{,}99$ относительная ошибка не превысила $6\cdot10^{-11}$.

{\singlespacing\small
\begin{Verbatim}
import numpy as np
from klein_surfaces import hyperbolic_sphere, horosphere

rng = np.random.default_rng(1)

def points(mesh, eps):
    k = np.stack([m.ravel() for m in mesh], 1)
    return k[1 - (k * k).sum(1) > eps]

def unit(v):
    return v / np.linalg.norm(v)

def distance(k, c):
    return np.arccosh((1 - k @ c) /
                      np.sqrt((1 - (k * k).sum(1)) * (1 - c @ c)))

def horo_const(k, omega):
    return (1 - k @ omega) / np.sqrt(1 - (k * k).sum(1))

worst_sphere = 0
for _ in range(200):
    c = unit(rng.normal(size=3)) * rng.uniform(0, 0.95)
    rho = rng.uniform(0.05, 2.5)
    k = points(hyperbolic_sphere(c, rho), 1e-12)
    worst_sphere = max(worst_sphere, np.abs(distance(k, c) - rho).max())

worst_horo = 0
for _ in range(200):
    omega = unit(rng.normal(size=3))
    r = rng.uniform(0.05, 0.95)
    k = points(horosphere(omega, r), 1e-6)
    expected = r / np.sqrt(1 - r * r)
    worst_horo = max(worst_horo,
                     np.abs(horo_const(k, omega) - expected).max())
print(worst_sphere, worst_horo)
\end{Verbatim}
}
